% ===================================================================
% preamble.tex - Praeambel des Vollberichts GEA Group
% Uebernommen aus Business_Analysen/Dual_Champions und auf einen
% einzelnen Titel zurueckgebaut: keine Teile, keine Paartabellen.
% ===================================================================
\usepackage[T1]{fontenc}
\usepackage[utf8]{inputenc}
\usepackage{lmodern}
\usepackage[ngerman]{babel}
\usepackage{csquotes}
% patch ohne "footnote": microtype kann seinen Fussnotenpatch nicht anlegen,
% wenn footmisc (para) den Fussnotenapparat ersetzt.
\usepackage[patch={item,toc,eqnum}]{microtype}
\usepackage{xcolor}
\usepackage{amsmath}
\usepackage{amssymb}
\usepackage{booktabs}
\usepackage{longtable}
\usepackage{array}
\usepackage{tabularx}
\usepackage{siunitx}
\usepackage{graphicx}
\usepackage{tikz}
\usepackage{placeins}
\usepackage{needspace}
\usepackage{pgfplots}
\usepackage[hidelinks]{hyperref}
\usepackage{url}

\areaset[current]{463.07pt}{636.6pt}

% Schusterjungen und Hurenkinder verbieten (im deutschen Satz Fehler).
\clubpenalty=10000
\widowpenalty=10000
\displaywidowpenalty=10000

% Fussnotenapparat flach halten: alle Fussnoten einer Seite in einem
% durchlaufenden Absatz. Das zweispaltige dblfnote laeuft mit pgfplots und
% hyperref in eine Endlosschleife der Ausgaberoutine.
\usepackage[para]{footmisc}

\usepackage[automark]{scrlayer-scrpage}
\pagestyle{scrheadings}
\renewcommand{\sectionmark}[1]{\markright{\thesection\quad #1}}
\clearpairofpagestyles
\setkomafont{pagehead}{\normalfont\small}
\ihead{\rightmark}
\ohead{GEA Group}
\KOMAoptions{headsepline=true}
\cfoot*{\pagemark}

\setcounter{tocdepth}{2}

\pgfplotsset{compat=1.18}
\usepgfplotslibrary{dateplot}
\pgfplotsset{/pgf/number format/use comma, /pgf/number format/1000 sep={}}

\definecolor{kursfarbe}{RGB}{25,55,120}
\definecolor{auftragsfarbe}{RGB}{0,120,60}

% --- Zahlenformat: deutsches Dezimalkomma ---------------------------
\sisetup{
  output-decimal-marker = {,},
  group-separator       = {.},
  group-minimum-digits  = 4,
  % Ohne group-digits=integer gruppiert siunitx auch die Nachkommastellen:
  % aus 43.4679 wurde im Satz "43,467.9".
  group-digits          = integer,
  detect-weight         = true,
  detect-family         = true
}

% --- Zitierweise ------------------------------------------------------
% \Q{Schluessel}{Seite}: Fussnote mit Kurztitel (Sprung ins Quellenverzeichnis)
% und GEDRUCKTER Seite. \QW{Titel}{Schluessel}{Abrufdatum}: Webquelle; die
% Adresse steht nur im Quellenverzeichnis, damit jede URL genau einmal vorkommt.
\definecolor{quellenlink}{RGB}{20,60,120}
\newcommand{\quellensprung}[2]{\hyperref[#1]{\textcolor{quellenlink}{#2}}}
\newcommand{\kurzGeaGB}{GEA Group, Annual Report 2025}
\newcommand{\kurzGeaGBvier}{GEA Group, Annual Report 2024}
\newcommand{\kurzGeaGBdrei}{GEA Group, Annual Report 2023}
\newcommand{\kurzGeaGBzwei}{GEA Group, Annual Report 2022}
\newcommand{\kurzGeaGBeins}{GEA Group, Annual Report 2021}
\newcommand{\kurzGeaGBnull}{GEA Group, Annual Report 2020}
\newcommand{\kurzGeaHJ}{GEA Group, Half-Year Financial Report 2026}
\newcommand{\Q}[2]{\footnote{\quellensprung{quelle:#1}{\csname kurz#1\endcsname}, S.~#2.}}
\newcommand{\QW}[3]{\footnote{\quellensprung{#2}{#1} (abgerufen am #3).}}

% --- Einheiten: Makros aus data/kennzahlen.tex sind bereits deutsch gesetzt
\newcommand{\Mio}[2]{#1\,Mio.\,#2}
\newcommand{\je}[2]{#1\,#2}
\newcommand{\Pz}[1]{#1\,\%}

% --- Drei Farben, drei Verantwortlichkeiten --------------------------
% Schwarz: belegte Fakten. Blau: Einschaetzung des Verfassers. Rot: Votum.
\definecolor{vdrot}{RGB}{155,25,30}
\newcommand{\vd}[1]{{\color{vdrot}#1}}
\definecolor{bkblau}{RGB}{20,60,150}
\newcommand{\bk}[1]{{\color{bkblau}#1}}

% --- Generierte Tabellenkoerper ohne angehaengtes \relax --------------
\makeatletter
\newcommand{\tabellenkoerper}[1]{\@@input #1.tex }
\makeatother

\newcommand{\annahme}[1]{%
  \par\medskip\noindent\fbox{\parbox{0.97\linewidth}{\small\textbf{Annahme:} #1}}\par\medskip}

% =====================================================================
% Tabellen. Jeder Koerper kommt aus data/ und wird von scripts/rechnung_gea.py
% erzeugt; im Dokument steht keine Zahl von Hand.
% =====================================================================
\newcommand{\ueberblicktabelle}{%
\begin{table}[htbp]\centering\small
\caption{Kennzahlen 2025 gegen 2024, wie auf der Kennzahlenseite gedruckt.
Ver\"anderung in Einheiten der Zeile (Prozentpunkte bei Quoten).}\label{tab:ueberblick}
\begin{tabular}{@{}lrrr@{}}\toprule
Kennzahl (Mio.~EUR, wo nicht anders angegeben) & 2025 & 2024 & Ver\"anderung\\\midrule
\tabellenkoerper{data/ueberblick}
\bottomrule\end{tabular}\end{table}}

\newcommand{\reihentabelle}{%
\begin{table}[htbp]\centering\small
\caption{Mehrjahresreihe 2019--2025 (Mio.~EUR; Marge und ROCE in \%, Ergebnis je Aktie in
EUR). \enquote{Freier Zufluss} ist die eigene Gr\"o{\ss}e dieses Berichts: operativer
Mittelzufluss der fortgef\"uhrten T\"atigkeit + Investitionen + Leasingtilgung; sie ist
nicht der von GEA ausgewiesene freie Cashflow (Methode S.~\pageref{sec:methode}).
Beleg: Gesch\"aftsbericht und gedruckte Seite der Kapitalflussrechnung.}\label{tab:reihe}
\begin{tabular}{@{}lrrrrrrrl@{}}\toprule
Jahr & Auftragseingang & Umsatz & EBITDA vR & Marge & ROCE & Freier Zufluss & EPS verw. & Beleg\\\midrule
\tabellenkoerper{data/reihe}
\bottomrule\end{tabular}\end{table}}

\newcommand{\prognosetabelle}{%
\begin{table}[htbp]\centering\small
\caption{Prognosetreue 2021--2025. Gelesen ist die Prognose \emph{wie zuerst gegeben},
nicht die unterj\"ahrig angepasste. Die Spalte \enquote{Ergebnis} stammt aus der eigenen
Reihe (Tabelle~\ref{tab:reihe}), nicht aus der Ergebnisspalte des Berichts; die Einordnung
rechnet \texttt{scripts/rechnung\_gea.py}.}\label{tab:prognose}
\begin{tabular}{@{}llllll@{}}\toprule
Jahr & Gr\"o{\ss}e & Prognose (zuerst) & Ergebnis & Lage & Seite\\\midrule
\tabellenkoerper{data/prognose}
\bottomrule\end{tabular}\end{table}}

\newcommand{\dividendentabelle}{%
\begin{table}[htbp]\centering\small
\caption{Dividende je Aktie, aus der Kapitalflussrechnung zur\"uckgerechnet: gezahlte
Dividende des Jahres geteilt durch die verw\"asserte Aktienzahl desselben Jahres. Die
Zahlung eines Jahres betrifft das Vorjahresergebnis.}\label{tab:dividende}
\begin{tabular}{@{}lrrrr@{}}\toprule
Jahr & Zahlung (Mio.~EUR) & Aktien (Mio.) & Je Aktie (EUR) & Schritt (EUR)\\\midrule
\tabellenkoerper{data/dividende}
\bottomrule\end{tabular}\end{table}}

\newcommand{\kurstabelle}{%
\begin{table}[htbp]\centering\small
\caption{Auftragseingang gegen Kurs: beide als Ver\"anderung zum Vorjahr.
Schlusskurse zum Jahresende, Yahoo Finance.}\label{tab:kurs}
\begin{tabular}{@{}lrrrr@{}}\toprule
Jahr & Auftragseingang (Mio.~EUR) & Ver\"and. in \% & Kurs Jahresende (EUR) & Ver\"and. in \%\\\midrule
\tabellenkoerper{data/kurs}
\bottomrule\end{tabular}\end{table}}

\newcommand{\optionentabelle}{%
\begin{table}[htbp]\centering\small
\caption{Interner Zinsfu{\ss} der drei Handlungsoptionen auf zehn Jahre, je Szenario des
freien Zuflusses. Option~C legt denselben Betrag zur H\"urde an; ihre Zeile ist die
Messlatte und rechnerisch szenariounabh\"angig.}\label{tab:optionen}
\begin{tabular}{@{}lrrrr@{}}\toprule
Option & Einstiegskurs (EUR) & pessimistisch & Basis & optimistisch\\\midrule
\tabellenkoerper{data/optionen}
\bottomrule\end{tabular}\end{table}}

\newcommand{\schwellentabelle}{%
\begin{table}[htbp]\centering\small
\caption{Schwellenkurs je Aktie (EUR): Kurs, zu dem der interne Zinsfu{\ss} die H\"urde von
\Pz{\Huerde} erreicht; letzte Spalte bei \Pz{\HuerdeLang}. Zufluss konstant, Endwert =
Buchwert des Eigenkapitals 2025.}\label{tab:schwelle}
\begin{tabular}{@{}lrrrrr@{}}\toprule
Szenario & Zufluss (Mio.~EUR) & 5 Jahre & 10 Jahre & 15 Jahre & 10 J., \Pz{\HuerdeLang}\\\midrule
\tabellenkoerper{data/schwelle}
\bottomrule\end{tabular}\end{table}}

\newcommand{\umkehrtabellen}{%
\begin{table}[htbp]\centering\small
\caption{Umkehrungen zum heutigen Kurs bei \Pz{\Huerde}: erforderliches Wachstum des freien
Zuflusses p.\,a.\ (oben) und erforderlicher Endwert als Vielfaches des Buchwerts
(unten).}\label{tab:umkehr}
\begin{tabular}{@{}lrrr@{}}\toprule
Szenario & 5 Jahre & 10 Jahre & 15 Jahre\\\midrule
\tabellenkoerper{data/wachstum}
\midrule
\tabellenkoerper{data/endwert}
\bottomrule\end{tabular}\end{table}}

\newcommand{\kgvtabelle}{%
\begin{table}[htbp]\centering\small
\caption{Kurs-Gewinn-Verh\"altnis zum Jahresende: Schlusskurs (Yahoo Finance) durch
verw\"assertes Ergebnis je Aktie (Tabelle~\ref{tab:reihe}); nur Jahre mit positivem
Ergebnis.}\label{tab:kgv}
\begin{tabular}{@{}lrrr@{}}\toprule
Jahr & Kurs Jahresende (EUR) & EPS verw. (EUR) & KGV\\\midrule
\tabellenkoerper{data/kgv}
\bottomrule\end{tabular}\end{table}}
